\documentclass[11pt]{article}

\usepackage[margin=0.78in]{geometry}
\usepackage[T1]{fontenc}
\usepackage{lmodern}
\usepackage{microtype}
\usepackage{amsmath,amssymb,mathtools}
\usepackage{booktabs,tabularx,array}
\usepackage{enumitem}
\usepackage[dvipsnames]{xcolor}
\usepackage{natbib}
\usepackage{hyperref}
\usepackage[nameinlink,noabbrev]{cleveref}

\newcolumntype{Y}{>{\raggedright\arraybackslash}X}
\newcolumntype{L}[1]{>{\raggedright\arraybackslash}p{#1}}
\setlength{\tabcolsep}{4.5pt}
\renewcommand{\arraystretch}{1.12}
\setlist{nosep,leftmargin=1.5em}

\newcommand{\utask}{U_{\mathrm{task}}}
\newcommand{\upers}{U_{\mathrm{pers}}}
\newcommand{\usafe}{U_{\mathrm{safe}}}
\newcommand{\cost}{C}
\newcommand{\history}{H_u^{(0)}}
\newcommand{\stream}{D_u^{(t)}}
\newcommand{\checkpoint}{c_u^{(t)}}

\newenvironment{settingpointer}
  {\begin{quote}\small\noindent\textbf{Prior-work setting.}\ }
  {\end{quote}}

\hypersetup{
  colorlinks=true,
  linkcolor=MidnightBlue,
  citecolor=ForestGreen,
  urlcolor=BrickRed,
  pdftitle={Problem Landscape and Success Criteria for One-User Offline-to-Online Personalized Policies}
}

\title{\bfseries Problem Landscape and Success Criteria for\\
One-User Offline-to-Online Personalized Policies}
\author{Research Position and Evaluation Contract}
\date{\today}

\begin{document}
\maketitle

\section{Setting}
\label{sec:setting}

\paragraph{Evidence base.}
Two project PRISMA runs yielded 6{,}663 raw records and 119 full-text reviews.
We cite primary papers for their methods and evaluation settings, not as
evidence that our policy works.

\paragraph{One user is the unit of personalization.}
For user $u$, the offline history
\[
  \history
  =
  \{\text{thoughts, conversations, corrections, preferences,
  task states, actions, tools, and outcomes}\}
\]
produces an initial checkpoint $c_u^{(0)}$.  After deployment, interaction
batch $\stream$ is generated by the active checkpoint $\checkpoint$.
The update is
\begin{equation}
  c_u^{(t+1)}
  =
  \mathcal A\!\left(c_u^{(t)},\history,D_u^{(1:t)}\right).
  \label{eq:update}
\end{equation}
Each record stores its user, time, task, checkpoint, trace, feedback source,
and training eligibility.  A shared prior is allowed; raw histories remain
separate.

\paragraph{A policy update is required.}
Retrieval, profiles, memory, and decoding control are baselines, not the
learned object.  OPPU trains a private user adapter
\citep{tan2024oppu}; P-RLHF trains a user-conditioned policy
\citep{li2024personalized}; and SPRInG updates a user adapter over time
\citep{kim2026spring}.  CoSteer is tuning-free
\citep{lv2025costeer}, while FABLE learns an external execution policy around
a frozen agent \citep{jin2026fable}.  Neither CoSteer nor FABLE updates a
deployed user's language-model policy weights.

\paragraph{Objective.}
Improve personalization without degrading task success or safety:
\begin{align}
  \max_{\mathcal A}\quad&
  \mathbb E[\upers(c_u^{(t+1)})] \nonumber\\
  \text{s.t.}\quad&
  \mathbb E[\utask(c_u^{(t+1)})-\utask(c_u^{(t)})]
    \ge -\delta_{\mathrm{task}}, \nonumber\\
  &
  \mathbb E[\usafe(c_u^{(t+1)})-\usafe(c_u^{(t)})]
    \ge -\delta_{\mathrm{safe}}, \nonumber\\
  &
  \cost_{\mathrm{train}},\cost_{\mathrm{serve}},
  \cost_{\mathrm{feedback}}
  \ \text{remain within declared budgets}.
  \label{eq:constrained-objective}
\end{align}
\paragraph{Exact scope.}
The target is one persistent user policy, initialized from that user's
history, updated from later same-user interactions, used for agent tasks, and
trained through language-model parameter updates.  We call an update
\emph{online} only when post-deployment evidence changes model parameters;
fresh rollouts within one training job do not qualify.  No cited method covers
all five properties.

\section{Related Problems and Setting Gaps}
\label{sec:problems}

\subsection{Learning the user}

\paragraph{P1: identify preferences from sparse, mixed evidence.}
A user supplies requests, choices, edits, and task outcomes, not one uniform
label.  Population training can blur minority preferences
\citep{li2024personalized,jang2023personalizedsoups}.  CoPL and PROPER share
structure across users \citep{choi2025copl,zhang2025proper}; FSPO adapts from
few examples \citep{singh2025fspo}.  We instead use one user's full history.

\begin{settingpointer}
\textbf{One user:} represented, but several methods train across users;
offline initialization: yes; post-deployment update: no; personalized agent:
no; LLM-policy training: P-RLHF, PROPER, and FSPO, but not CoPL.
\textbf{Difference:} fixed-data preference learning, not deployed agent updates.
\end{settingpointer}

\paragraph{P2: separate persistent preference from transient context.}
A request may express a lasting preference, a task constraint, or noise.
SPRInG filters for drift before updating a user adapter
\citep{kim2026spring}.  Agent traces add tool and reasoning failures that are
not preference evidence.

\begin{settingpointer}
\textbf{One user:} yes; offline initialization: yes;
post-deployment update: no; personalized agent: no; LLM-policy training: yes.
\textbf{Difference:} SPRInG replays a fixed chronological corpus rather than live
agent trajectories.
\end{settingpointer}

\paragraph{P3: handle cold start without losing ownership.}
Collaborative and federated models improve sparse-user adaptation
\citep{choi2025copl,zhang2025proper,zhu2026fedpdpo}.  Our setting may use a
shared prior, but histories, updates, and evaluation remain user-specific.

\begin{settingpointer}
\textbf{One user:} represented through shared training; offline
initialization: yes; post-deployment update: no; personalized agent: no;
LLM-policy training: PROPER and FedPDPO, while CoPL trains a reward model.
\textbf{Difference:} these methods pool users offline; our known-user setting
does not require cold start.
\end{settingpointer}

\subsection{Moving from offline history to online learning}

\paragraph{P4: build an offline initializer with online headroom.}
The offline checkpoint must complete tasks yet remain adaptable.  OPPU
provides a per-user initializer \citep{tan2024oppu}; WebRL shows the value of
warm starts for sparse-reward agents \citep{qi2025webrl}.

\begin{settingpointer}
\textbf{One user:} OPPU only; offline initialization: yes;
post-deployment update: no; personalized agent: WebRL only; LLM-policy
training: yes.  \textbf{Difference:} neither paper combines these properties.
\end{settingpointer}

\paragraph{P5: cross the offline--online boundary without a performance dip.}
Offline records come from older policies.  Synchronization
\citep{lanchantin2025bridging}, hybrid preference learning
\citep{bose2024hybrid}, and policy expansion \citep{zhang2023pex} address this
handoff.  Our handoff must also preserve user, feedback, and checkpoint
lineage.

\begin{settingpointer}
\textbf{One user:} no; offline initialization: yes;
post-deployment update: no; personalized agent: no; LLM-policy training:
Bridging only.  \textbf{Difference:} their ``online'' data are collected during
training, not from a deployed user.
\end{settingpointer}

\paragraph{P6: adapt to drift without forgetting or chasing noise.}
SPRInG selects likely preference shifts and retains residual memory
\citep{kim2026spring}.  Stable Asynchrony handles stale trajectories from
lagged policies \citep{huang2026stable}.  We need both without forgetting
stable preferences.

\begin{settingpointer}
\textbf{One user:} SPRInG only; offline initialization: SPRInG only;
post-deployment update: no; personalized agent: no; LLM-policy training: yes.
\textbf{Difference:} neither joins preference drift and stale agent trajectories
after deployment.
\end{settingpointer}

\paragraph{P7: learn quickly while minimizing user burden.}
User feedback is expensive.  Hybrid preference optimization trades offline
data against online queries \citep{bose2024hybrid}.  FABLE learns a pure-online
per-user execution policy \citep{jin2026fable}.  We measure interactions
needed to reach a fixed future-task utility.

\begin{settingpointer}
\textbf{One user:} FABLE only; offline initialization: hybrid preference
optimization only; post-deployment update: FABLE only; personalized agent:
FABLE only; LLM-policy training: not FABLE.  \textbf{Difference:} FABLE updates
an external controller, not an initialized LLM policy.
\end{settingpointer}

\subsection{Learning an agent policy}

\paragraph{P8: preserve task completion while personalizing execution.}
Personalized agents choose tools, memory, clarifications, and response style
while completing the task.  PersonalWAB studies web actions
\citep{cai2025personalwab}, FingerTip 20K studies mobile execution
\citep{yang2026fingertip}, and FABLE adapts execution around a frozen agent
\citep{jin2026fable}.

\begin{settingpointer}
\textbf{One user:} persistent policy only in FABLE; offline initialization: no
persistent per-user checkpoint; post-deployment update: FABLE only;
personalized agent: yes; LLM-policy training: not across both phases.
\textbf{Difference:} none trains one personal LLM across both phases.
\end{settingpointer}

\paragraph{P9: assign credit from long, partially observed trajectories.}
A correction may follow many reasoning and tool steps.  WebRL learns from
sparse terminal feedback \citep{qi2025webrl}; we must also identify which
steps reveal the user's preference.

\begin{settingpointer}
\textbf{One user:} no; offline initialization: yes;
post-deployment update: no; personalized agent: no; LLM-policy training: yes.
\textbf{Difference:} WebRL's credit is task-conditioned, not user-conditioned.
\end{settingpointer}

\subsection{Obtaining credible and affordable evidence}

\paragraph{P10: use simulators without training to simulator artifacts.}
UserRL, UserVille, ProPerSim, and CollabLLM train or evaluate agents with
simulated users
\citep{qian2025userrl,sun2025userville,kim2025propersim,wu2025collabllm}.
PersonalLLM creates synthetic preference models \citep{zollo2024personalllm}.
Simulators reduce cost but can impose their own style and reward biases
\citep{rezk2025rmcrisis}; real-user evidence remains necessary.

\begin{settingpointer}
\textbf{One user:} simulated personas, not a persistent target policy; offline
initialization: method-dependent; post-deployment update: no; personalized
agent: several methods; LLM-policy training: method-dependent.
\textbf{Difference:} simulated users cannot establish improvement for one
deployed real user.
\end{settingpointer}

\paragraph{P11: make per-user deployment cheaper than repeated large-model use.}
Per-user adapters and online training add storage and compute.  OPPU uses
lightweight adapters \citep{tan2024oppu}, CoSteer steers a cloud model from a
local model \citep{lv2025costeer}, and FABLE learns a small external policy
\citep{jin2026fable}.  We test whether a compact user policy can match a
larger generic model at lower serving cost.

\begin{settingpointer}
\textbf{One user:} yes; offline initialization: OPPU only;
post-deployment update: FABLE only; personalized agent: FABLE only; LLM-policy
training: OPPU only.  \textbf{Difference:} CoSteer is tuning-free and FABLE trains an
external policy; none covers both phases.
\end{settingpointer}

\paragraph{P12: evaluate future behavior rather than history recall.}
LongLaMP provides chronological long-form tasks \citep{kumar2024longlamp};
PrefEval tests preference following \citep{zhao2025prefeval}; PERMA adds
evolving preferences and tasks \citep{liu2026perma}; and PRISM supplies
individual feedback \citep{kirk2024prism}.  Our endpoint is behavior on
untouched future tasks.

\begin{settingpointer}
\textbf{One user:} represented in the data; offline initialization: not
applicable; post-deployment update: no; personalized agent: PERMA only;
LLM-policy training: no.  \textbf{Difference:} these are evaluation resources,
not update algorithms.
\end{settingpointer}

\section{Claim}
\label{sec:claim}

\begin{quote}
\textbf{A compact, user-owned language-model policy can be initialized from
one person's heterogeneous historical interactions and then improved from
that person's later agent-task evidence through real parameter updates, while
preserving task success and safety.}
\end{quote}

This requires an offline personal checkpoint, later parameter updates, and
paired evaluation of personalization and task success.  The method
contribution is the update rule connecting history, fresh trajectories,
feedback, replay, and checkpoint activation.

\subsection{Additional claims}

Report these claims only when their tests pass:
\begin{description}
\item[Fast adaptation.] Fewer real interactions reach the same valid utility.
\item[Trajectory advantage.] Tool actions and verified outcomes improve
  future execution beyond text-only history.
\item[Small versus large.] A compact personalized policy beats a larger
  generic model while preserving task success and lowering cost.
\item[Simulation benefit.] Simulator training reduces real-user interactions
  and survives held-out human evaluation.
\item[Privacy.] The implementation follows its declared data boundary without
  exposing another user's raw records.
\end{description}

\section{Success Criteria}
\label{sec:success}

\subsection{Hard validity gates}

Before comparing scores, every reported update must pass:
\begin{enumerate}
  \item the parent and candidate checkpoint hashes differ;
  \item the candidate loads in a fresh worker and actually generates the
  evaluated trajectories;
  \item offline, online-adaptation, validation, and future-test intervals are
  chronological and disjoint;
  \item every training record has user, task, behavior-policy, checkpoint,
  timestamp, and feedback provenance;
  \item failed runs, users, and rollbacks remain in the report;
  \item the simulator, reward model, or LLM judge has not seen the future-test
  target or its answer.
\end{enumerate}
Failure of any item invalidates the run rather than lowering its score.

\subsection{Primary statistical gate}

Let $b$ be the strongest eligible baseline on the validation set.  For each
user-period cluster, evaluate the candidate and $b$ on identical future
tasks.  Define
\begin{align}
  \Delta_{\mathrm{pers}}
    &= \upers(c_{\mathrm{candidate}})-\upers(c_b),\\
  \Delta_{\mathrm{task}}
    &= \utask(c_{\mathrm{candidate}})-\utask(c_b).
\end{align}
The main claim passes only when
\begin{equation}
  \operatorname{LCB}_{95\%}(\Delta_{\mathrm{pers}})>0
  \quad\text{and}\quad
  \operatorname{LCB}_{95\%}(\Delta_{\mathrm{task}})
  \ge -\delta_{\mathrm{task}}.
  \label{eq:primary-gate}
\end{equation}
The minimum meaningful personalization effect and
$\delta_{\mathrm{task}}$ must be fixed before final evaluation.  Confidence
intervals are clustered by user or task lineage, not by treating correlated
turns as independent samples.

\begin{table}[t]
  \centering
  \caption{Success criteria.  The first four rows test the main claim.}
  \label{tab:success}
  \small
  \begin{tabularx}{\textwidth}{@{}L{0.18\textwidth}L{0.27\textwidth}Y@{}}
    \toprule
    Claim & Measurement & Pass condition \\
    \midrule
    Personalized improvement &
    Paired future-task $\Delta_{\mathrm{pers}}$ against the strongest baseline. &
    Lower 95\% bound is above zero and the point estimate exceeds the
    preregistered meaningful effect. \\
    Task preservation &
    Paired task success or environment reward. &
    Lower 95\% bound exceeds $-\delta_{\mathrm{task}}$. \\
    Value of online learning &
    Candidate versus frozen offline checkpoint and naive cumulative update. &
    Candidate passes \cref{eq:primary-gate} against both. \\
    Real policy change &
    Parameter count, hashes, load receipt, and changed outputs. &
    Nonzero finite update; candidate is loaded for evaluation. \\
    Fast adaptation &
    Area under the future-utility learning curve and interactions-to-target. &
    Higher clustered 95\% interval for area, or fewer interactions to the
    same preregistered target. \\
    Retention &
    Backward transfer on earlier stable preferences and tasks. &
    No task-margin violation; less forgetting than naive continual training. \\
    Drift recovery &
    Utility after a controlled persistent preference change. &
    Faster recovery without elevated updates on transient-noise controls. \\
    Safety &
    Severe violations and policy regressions under a fixed panel. &
    No increase in severe events; lower-severity change stays within its
    preregistered margin. \\
    Small versus large &
    Compact personalized model versus large generic, RAG, and profile-enabled
    models under equal information. &
    Personalization superiority, task non-inferiority, and strictly lower
    serving cost all hold. \\
    Simulator benefit &
    Human--simulator agreement, calibration, and method ranking. &
    Simulator-trained gain survives held-out human evaluation; simulation is
    never the only final judge. \\
    Efficiency &
    Training GPU-hours, inference latency, memory, storage, tokens, API cost,
    and user-feedback count. &
    Candidate lies on the quality--cost Pareto frontier under matched budgets. \\
    Ownership/privacy &
    Raw-record flow, communication bytes, cross-user leakage probes. &
    Declared boundary is met with no unauthorized raw-history transfer. \\
    \bottomrule
  \end{tabularx}
\end{table}

\subsection{Large-model comparison}

The large-model comparison must include:
\begin{enumerate}
  \item a large generic model without personal history;
  \item the same model with the full allowed history;
  \item strong RAG or profile prompting;
  \item a decoding-time method such as CoSteer when executable
  \citep{lv2025costeer};
  \item the compact offline-only and compact online-adapted policies.
\end{enumerate}
Match prompts, tools, history, and task budgets.  A compact model wins only if
it improves personalization, preserves task success, and costs less.

\subsection{Adaptation speed}

For threshold $\tau$, measure real interactions rather than epochs:
\begin{equation}
  T_{\tau}
  =
  \min\{n:\upers(c_u^{(n)})\ge\tau
  \ \text{and task/safety gates pass}\}.
\end{equation}
Also report learning-curve area, final utility, clarifications, correction
time, and rollbacks.

\subsection{Simulator evidence}

Test:
\begin{enumerate}
  \item \textbf{Preference fidelity:} agreement and calibration on held-out
  real-user comparisons.
  \item \textbf{Method fidelity:} rank correlation between methods under the
  simulator and under real-user evaluation.
  \item \textbf{Transfer value:} fewer real interactions are needed after
  simulator pretraining than without it.
\end{enumerate}
Simulator reward alone does not support a real-user claim.

\section{Experimental Comparison Set}
\label{sec:baselines}

\begin{table}[t]
  \centering
  \caption{Minimum baseline families.}
  \label{tab:baselines}
  \small
  \begin{tabularx}{\textwidth}{@{}L{0.23\textwidth}YY@{}}
    \toprule
    Family & Purpose & Representative implementation \\
    \midrule
    No personalization &
    Measures the value of any user information. &
    Same-size base and large API model. \\
    Prompt/profile/RAG &
    Tests whether parameter updates are necessary. &
    Full-history prompting, retrieval, and generated profile. \\
    Offline per-user PEFT &
    Tests the value of the online phase. &
    OPPU \citep{tan2024oppu}. \\
    Shared personalized policy &
    Tests one-user ownership against a population model. &
    P-RLHF, CoPL, or PROPER
    \citep{li2024personalized,choi2025copl,zhang2025proper}. \\
    Continual per-user update &
    Tests the proposed transition and replay rule. &
    SPRInG and naive cumulative SFT \citep{kim2026spring}. \\
    Online-only &
    Measures the value of offline initialization. &
    Same optimizer initialized from the generic checkpoint. \\
    Offline then naive online &
    Isolates the new coupling mechanism. &
    Sequential SFT/DPO/RL with no special replay or activation rule. \\
    Weight-free online control &
    Tests whether policy training is worth its cost. &
    FABLE and CoSteer \citep{jin2026fable,lv2025costeer}. \\
    Large-model teacher &
    Tests compact-model competitiveness. &
    The CLI/API model that produced the offline trajectories. \\
    \bottomrule
  \end{tabularx}
\end{table}

\paragraph{Evaluation substrates.}
Use a chronological stream such as LongLaMP \citep{kumar2024longlamp} and an
executable agent benchmark such as PersonalWAB or FingerTip 20K
\citep{cai2025personalwab,yang2026fingertip}.  PrefEval and PERMA add
long-horizon tests \citep{zhao2025prefeval,liu2026perma}.  Final results need
held-out human feedback or frozen real-user records and two distinct
substrates.

\paragraph{Replication.}
Repeat the protocol across independent users.  Report per-user results, the
median, worst decile, and cluster-aware intervals, with at least three seeds
for feasibility and five for final results.

\section{Claims We Must Not Make}
\label{sec:nonclaims}

\begin{itemize}
  \item Chronological offline replay is not post-deployment online learning.
  \item Better retrieval or profile recall is not evidence of a better policy.
  \item A simulator-only win is not evidence of real-user improvement.
  \item A changed loss value is not evidence that the candidate checkpoint was
  loaded.
  \item A personalized reward-model gain is not automatically a policy or task
  gain \citep{rezk2025rmcrisis}.
  \item Local adapters do not by themselves prove privacy.
  \item Beating a large model on style does not mean beating it on task
  capability.
  \item Mean improvement cannot hide severe failures for a subset of users.
\end{itemize}

\bibliographystyle{plainnat}
\bibliography{personalization_problem_landscape}

\end{document}
